%% document class elsarticle, preprint option
\documentclass[preprint,12pt]{elsarticle}
%% Use the option review to obtain double line spacing

%%%%%%%%%%%%%%%%%%%%%%%%
%% installed packages %%
%%%%%%%%%%%%%%%%%%%%%%%%

\usepackage{amssymb}
\usepackage{amsmath}
\usepackage{booktabs} % idk what this does
\usepackage{siunitx} % support for SI units
\usepackage{tcolorbox} % framework boxes and worked example
\tcbuselibrary{breakable,skins} % idk what this does
\usepackage{algorithm} % formal algorithm environment
\usepackage{algpseudocode} % formal algorithm environment
\usepackage{multirow}
\usepackage[section]{placeins} % %% placeins: prevent floats from drifting past section/subsection boundaries

\makeatletter %% Also barrier at subsections:
\AtBeginDocument{%
  \expandafter\renewcommand\expandafter\subsection\expandafter
    {\expandafter\@fb@secFB\subsection}%
}
\makeatother

\usepackage{hyperref} % for \href links

%%%%%%%%%%%%%%%%%%%%
%% Document begin %%
%%%%%%%%%%%%%%%%%%%%

\journal{Solar Energy}

\begin{document}

\begin{frontmatter}

%%%%%%%%%%%%%%%%%%%%%%%%%%%%%%%%%%
%% Title, authors and addresses %%
%%%%%%%%%%%%%%%%%%%%%%%%%%%%%%%%%%

%% article title
\title{Estimation of Longwave All-sky Emissivity using Shortwave and Meteorological Measurements} 

%% authors and contributions 
\author[label1]{Akarsh Srivastava \fnref{fn_equal}}

\author[label1]{Jiedong Wang\fnref{fn_equal}} 

\author[label1]{Zoey Huang} 

\author[label1]{Clayson Briggs} 
 
\author[label1]{Carlos F.\,M. Coimbra \corref{correspond}}

\fntext[fn_equal]{These authors contributed equally to this work and share first authorship.}
\cortext[correspond]{Corresponding author: ccoimbra@ucsd.edu}

%% Author affiliation
\affiliation[label1]{
    organization={Department of Mechanical and Aerospace Engineering, University of California San Diego},
    city={La Jolla},
    postcode={92093-0411}, 
    state={California},
    country={USA} 
}

%%%%%%%%%%%%%%
%% Abstract %%
%%%%%%%%%%%%%%

\begin{abstract}
Downwelling longwave radiation (DLR) is a critical component of the surface energy budget, yet routine measurements remain sparse globally due to cost and difficulty of measurement. 
% Fewer than 5\% of global meterological stations have these DLR measurements using pyrgeometers, while pyranometers for shortwave radiation measurements are deployed at the majority of the WMO-registered stations.\citep{wang2013global}. 
% Furthermore BSRN network only has $\sim$60 sites globally that provide research-grade DLR measurements.\citep{driemel2018baseline}. 
% World Meterorological Organization's Global Observing System has over 11,000 surface stations but fewer than 200 measure DLR routinely. \citep{wang2013global}.
%This data gap introduces significant uncertainties in the global energy budget and complicates the evaluation of the climate model simulations\citep{wild2017new,stephens2012global}. 
This study builds on previous work in clear-sky DLR estimation, using a cloud factor $\gamma$ to quantify cloud influences on atmospheric emissivity. Three parametric models relating $\gamma$ to shortwave radiation measurements are calibrated with balanced sampling of 13 years (2010-2022) of high-precision data from seven SURFRAD stations across diverse geographies and climates. The model based on diffuse fraction ($k_d$) achieves low error (RMSE $=0.030$) across all sites and years when compared to similar models based on clearness indices ($k_t$). This model is validated using independent 2023–2024 SURFRAD data, achieving a station-averaged RMSE of 0.037, rRMSE of 4.6\%, and $R^2$ of 0.83, indicating the model's robustness. The model also remains accurate (rMBE $= 2.7 \% -6.1 \%$ and rRMSE $= 5.4 \% - 10.0 \%$) when using empirical correlations for diffuse fraction, extending its applicability beyond specialized networks. In total, the proposed approach enables more accurate DLR estimation under all-sky conditions, thereby enhancing surface energy budget analyses in data-sparse regions.

%This study addresses this data gap by parameterizing the all-sky DLR and effective sky emissivity as a function of widely available shortwave radiation measurements. An effective cloud fraction factor, notated $\gamma$, is introduced to quantify cloud influences on atmospheric emissivity. A simple power-law function of diffuse fraction ($k_d$) is found to correlate with $\gamma$ and is calibrated with balanced sampling of 13 years (2010-2022) of high-precision data from seven SURFRAD stations across diverse geographies and climates. The model achieves low error (RMSE $=0.030$) across all sites and years when compared to similar models based on clearness indices ($k_t$). This model is validated using independent 2023–2024 SURFRAD data, achieving a station-averaged RMSE of 0.037, rRMSE of 4.6\%, and $R^2$ of 0.83, indicating the model's robustness. The model also remains accurate (rMBE $= 2.7 \% -6.1 \%$ and rRMSE $= 5.4 \% - 10.0 \%$) when using empirical correlations for diffuse fraction, extending its applicability beyond specialized networks. In total, the proposed approach enables more accurate DLR estimation under all-sky conditions, thereby enhancing surface energy budget analyses in data-sparse regions.
\end{abstract}

%%Graphical abstract 
%\begin{graphicalabstract}
%\includegraphics{grabs}
%\end{graphicalabstract}

%% Research highlights
\begin{highlights}
\item Shortwave radiation measurements effectively parameterize cloud factor
\item Cloud factor enables longwave downwelling radiation estimation in all-sky conditions 
\item Parameterized model is robust across diverse geographies and climates
\end{highlights}

%% Keywords
\begin{keyword}
%% keywords here, in the form: keyword \sep keyword
Surface radiation budget \sep Downwelling longwave radiation \sep Atmospheric emissivity \sep Parametric modeling

%% PACS codes here, in the form: \PACS code \sep code

%% MSC codes here, in the form: \MSC code \sep code
%% or \MSC[2008] code \sep code (2000 is the default)

\end{keyword}

\end{frontmatter}

%% Add \usepackage{lineno} before \begin{document} and uncomment 
%% following line to enable line numbers
%% \linenumbers

%%%%%%%%%%%%%%%
%% main text %%
%%%%%%%%%%%%%%%

%% Use \section commands to start a section
\section{Introduction}
\label{sec1}
%% Labels are used to cross-reference an item using \ref command.

Downwelling longwave radiation (DLR) is a significant component of radiative heat flux incident on the surface, primarily driven by thermal emissions from the low atmosphere \cite{idso1969thermal,grobner2009effective,alados2012estimation}. 
Accurate values of DLR are important for several engineering applications, including the design of photovoltaic panels \cite{duffie2013solar}, as well as for weather forecasting and climate modeling \cite{trenberth2009earth,wild2017new}. % For example, during arctic summers, DLR serves as the dominant driver of sea ice melt in the net surface energy budget, and errors in its estimation can significantly affect climate model predictions of the extent of sea-ice in September \citep{luo2023uncertain}. 
While DLR can be directly measured using pyrgeometers, these instruments are expensive and challenging to calibrate, so many routine surface radiation measurement stations have historically focused on shortwave measurements \citep{gilgen1998means,wang2009global,stephens2012global,wang2013global}. Some advanced networks, such as NOAA's Surface Radiation Budget Network (SURFRAD) or the WMO's Baseline Surface Radiation Network (BSRN), provide reliable DLR data, but comprehensive global coverage remains limited, especially in developing regions. %This lack of routine global DLR measurements introduces uncertainties in the global surface energy budget, which in turn complicates meteorological and climate model simulations \citep{wild2017new,stephens2012global}. 
Consequently, substantial efforts have been directed toward developing simplified parameterized models which characterize the radiative properties of the sky as functions of readily available meteorological and radiometric measurements.  

% due to the high cost and technical challenges associated with the deployment of high-precision longwave radiometers, many routine surface radiation networks have historically focused on shortwave measurements \citep{stephens2012global}. Although advanced networks like SURFRAD provides reliable downwelling longwave radiation (DLR) data, comprehensive global coverage remains limited, especially in developing regions, The lack of routine DLR measurements introduces uncertainties in the global energy budget, which in turn complicates the evaluation of climate model simulations \citep{wild2017new}. Consequently, substantial efforts have been directed toward developing simplified parameterization methods based on readily measurable meteorological variables.

Significant progress has been made in predicting clear-sky emissivity, with numerous formulations available in the literature. These parameterization approaches rely on a combination of theoretical considerations and empirical relationships between the atmospheric emissivity and the moisture content and temperature of the air at the surface (usually measured $\SI{2}{\meter}$ above the surface, referred to as the screening height). Brunt \cite{brunt1932notes} proposed an early and widely used linear regression model that incorporates dew point temperature and near-surface air temperature. Authors such as Idso and Jackson \cite{idso1969thermal} and Brutsaert \cite{brutsaert1975derivable} later attempted to provide more universally applicable models by incorporating updated atmospheric radiation theory. More recently, Matsunobu et al. \cite{matsunobu2024effective} used a large dataset with high temporal resolution and measurement accuracy to optimize model fitting and include altitude corrections, thereby improving both model precision and applicability. 

However, accurate DLR estimates rely on all-sky emissivity predictions, where the complexity of cloud cover poses a major challenge to developing universally applicable models. The properties of clouds, such as amount, height, optical thickness, droplet size, and phase distribution, all influence the surface heat budget on broad temporal and spatial scales \citep{kiehl1997earth,li2017determination,kay2016recent}. %Due to the challenges of parameterizing these complex effects, a globally robust parameterization for all-sky emissivity remains elusive. 
One commonly employed strategy is to calculate all-sky emissivity by combining clear-sky emissivity with a cloud correction factor. Several groups \cite{arnfield1979evaluation,alados1995estimation,niemala2001comparison} have evaluated and compared parameterizations that use Boltz-type correction factors, which account for cloud type and cloud cover fraction. Several other groups \cite{alados2012estimation,crawford1999improved} have introduced cloud correction factors which use measurements of shortwave solar irradiance as a proxy for cloudiness. The present study extends the work of these latter groups, by introducing and comparing all-sky emissivity models based on several shortwave measurements from large datasets spanning the contiguous United States, offering a more widely applicable parameterization for all-sky emissivity calculations. 

\section{Background}
DLR is the thermal infrared energy emitted by atmospheric gases and clouds toward the Earth’s surface. This radiation spans wavelengths from approximately $\SI{4}{\micro\meter}$ to $\SI{100}{\micro\meter}$, with the majority of its energy concentrated in the $\SIrange{8}{14}{\micro\meter}$ spectral window \cite{kiehl1997earth}. In this longwave band, the atmosphere may be reasonably approximated as a gray body, and the effective emissivity of the sky $\varepsilon_{\mathrm{sky}}$ can be expressed in terms of DLR and the screening height temperature $T_a$ as \cite{brunt1932notes,matsunobu2024effective,berdahl2021retrospective}

\begin{equation} \label{sky_emissivity}
\varepsilon_{\mathrm{sky}} \, = \, \frac{\mathrm{DLR}}{\sigma\,T_a^4},
\end{equation}
where $\sigma$ is the Stefan-Boltzmann constant equal to $\SI{5.67e-8}{\watt\per\meter^{2}\kelvin}$. 

To estimate DLR, Eq.~\eqref{sky_emissivity} is typically rearranged, and measurements of atmospheric temperature and available data or correlations for emissivity are used. In the clear-sky case, a number of studies have already developed models for estimating emissivity, and the formula proposed by \cite{matsunobu2024effective} is adopted here due to its accuracy across different climate conditions. The expression for clear-sky emissivity is 

\begin{equation} \label{e_clear sky}
    \varepsilon_{\mathrm{clear\,sky}} =  0.600 + 1.652 \, \sqrt{p_w} + 0.150 \, \left[ \exp\left( -\frac{z}{8500} \right) - 1 \right],
\end{equation}
where $z$ is the site altitude in meters and $p_w$ is the local dimensionless partial pressure of water vapor at the screening height, defined as 
\begin{equation}
     p_{w} = \frac{P_{w}}{P_{0}}.
\end{equation}
Here, $P_0$ is the standard atmospheric pressure at sea level, taken as $\SI{101325}{\pascal}$, and $P_{w}$ is the actual partial pressure of water vapor. The value of $P_w$ can be found using the relative humidity $\phi$ and the saturation vapor pressure $P_s$ from
\begin{equation}
   P_w = P_s(T_a) \cdot \phi,
\end{equation}
and the saturation vapor pressure $P_s(T_a)$, in Pascals, can be expressed as a function of the air temperature $T_a$, in Kelvin, using the empirical equation \citep{alduchov1996improved}: 
\begin{equation}
   P_s(T_a) = 610.94 \, \exp \left( \frac{17.625 \,(T_a - 273.15)}{T_a - 30.11 } \right).
\end{equation}

In the all-sky case, accurate values of emissivity are challenging to obtain due to the complexities of cloud cover. Most methods used to find an all-sky emissivity value involve adding a corrective term to the longwave radiation estimated under clear-sky conditions \cite{alados2012estimation,arnfield1979evaluation,alados1995estimation,crawford1999improved}. Most recently, Coimbra \cite{coimbra2025energy} proposed an effective total emissivity model for the cloudy sky, incorporating an effective cloud factor $\gamma$:

\begin{equation} \label{e_cloudy sky}
\varepsilon_{\mathrm{sky}} = (1-\varepsilon_{\mathrm{clear\,sky}})\,  \gamma + \varepsilon_{\mathrm{clear\,sky}}.
\end{equation}
This cloud factor formulation normalizes the all-sky emissivity with respect to the clear-sky emissivity, taking on the value of $\gamma = 0$ when $\varepsilon_{\mathrm{sky}} = \varepsilon_{\mathrm{clear\,sky}}$ and the value of $\gamma = 1$ when $\varepsilon_{\mathrm{sky}} = 1$. Note that $\gamma$ can take on negative values in the case of a cloud layer being warmer than the ground, i.e. an atmospheric inversion layer. 

This study focuses on determining the effective cloud fraction factor $\gamma$ so that the all-sky emissivity may be estimated. To do this, the effects of cloud cover are quantified through its effects on three shortwave irradiance measurements: global horizontal irradiance (GHI), diffuse horizontal irradiance (DHI), and direct normal irradiance (DNI). Two quantities are commonly employed to compare these measurements in all-sky applications: 
\begin{enumerate}
    \item The clearness index $k_t$ is defined as the ratio of the measured irradiance to the theoretical clear-sky irradiance. This index can be computed using either GHI or DNI:
    \begin{equation} \label{kt_ghi}
     k_{t, \mathrm{GHI}} = \frac{\mathrm{GHI}_{\mathrm{measured}}}{\mathrm{GHI}_{\mathrm{ clear\, sky}}},
    \end{equation}
    \begin{equation} \label{kt_dni}
     k_{t,\mathrm{DNI}} = \frac{\mathrm{DNI}_{\mathrm{measured}}}{\mathrm{DNI}_{\mathrm{ clear\, sky}}}.
    \end{equation}

    \item The diffuse fraction $k_d$ is defined as the ratio of DHI to GHI:
    \begin{equation} \label{kd}
     k_d = \frac{\mathrm{DHI}_{\mathrm{measured}}}{\mathrm{GHI}_{\mathrm{measured}}}.
    \end{equation}
\end{enumerate}

Clear-sky GHI and DNI values are acquired from %the \texttt{location.get\_clearsky} function from the open-source pvlib library is employed. This function utilizes 
the Ineichen/Perez model \cite{ineichen2002new,ineichen2008broadband,ineichen2016validation}, which requires location-specific inputs such as latitude, longitude, altitude, and solar time. All necessary parameters can be obtained from the dataset used in this study without additional assumptions or external input.

The $k_t$ indices in Eq.s~\eqref{kt_ghi}-\eqref{kt_dni} indicate the degree of total irradiance attenuation caused by cloud cover and $k_d$ in Eq.~\eqref{kd} characterizes the changes in the spectral energy distribution induced by clouds. Together, these metrics provide a framework for quantifying the radiative effects of clouds. The effective cloud fraction factor can then be expressed as a function of $k_t$ and $k_d$:

\begin{equation} \label{fungamma}
  \gamma  = f\left( k_{t,\mathrm{GHI}}, \, k_{t,\mathrm{DNI}}, \, k_d\right).
\end{equation}

Based on Eq.~\eqref{fungamma}, we propose three parametric models to predict $\gamma$, and therefore the all-sky atmospheric emissivity and the downwelling longwave radiation, and evaluate them against each other using high-precision data and standard error metrics. These models establish a nonlinear relationship between the cloud fraction factor and shortwave radiation parameters, formulated as  

\begin{equation} \label{model_1} 
    \gamma = c_1\,k_d^{c_2}\,,
\end{equation}

\begin{equation} \label{model_2}
\gamma = c_1 \,k_{t,\mathrm{GHI}}^{c_2}\,,
\end{equation}

\begin{equation} \label{model_3}
 \gamma = c_1 \,k_{t,\mathrm{DNI}}^{c_2}.
\end{equation}
where $c_1$ and $c_2$ are fitting coefficients. 

Equation~\eqref{model_1} showcases Model 1, which relates the cloud fraction factor to the diffuse fraction $k_d$, Eq.~\eqref{model_2} shows Model 2, relating to the clear-sky index of global horizontal irradiance $k_{t,\mathrm{GHI}}$, and Eq.~\eqref{model_3} shows Model 3, based on the clear-sky index of direct normal irradiance $k_{t,\mathrm{DNI}}$. After obtaining optimal fitting results, the predicted values are compared to the observed data DLR and emissivity data. The model that exhibits the strongest correlation and lowest prediction error is then identified and analyzed. 

\section{Data Processing}

\subsection{Dataset Description}

In this study, the training data are sourced from the Surface Radiation Budget Network (SURFRAD), which includes seven radiation monitoring stations distributed across North America. Established in 1993 and managed by the National Oceanic and Atmospheric Administration (NOAA) Office of Global Programs, the network was designed to support long-term, continuous, and high-accuracy observations for studies on surface radiation balance, climate model validation, and satellite remote sensing data calibration. To ensure this, SURFRAD follows the measurement standards set by the Baseline Surface Radiation Network, sponsored by the World Climate Research Program of the World Meteorological Organization. Each station provides high-quality radiation and meteorological data, including global horizontal irradiance (GHI), diffuse horizontal irradiance (DHI), direct normal irradiance (DNI), downwelling longwave radiation (DLR), and relative humidity ($\phi$). %The achievable accuracy for shortwave radiation measurements is approximately 2\% for pyrheliometers and 5\% for pyranometers. 

The seven SURFRAD stations span a large range of geographies, altitudes, and climate conditions. This spatial and climatic diversity forms a comprehensive observation matrix that enhances the parametric models' generalization and global applicability. The geographic distribution of these stations is shown in Fig.~\ref{fig:surfrad_map}, and a summary of relevant parameters for each station is included in Table~\ref{tab:station-data}, such as geographic information and average values of meteorological measurements throughout the year. In particular, the cloud frequency in Table~\ref{tab:station-data} is defined as the percentage of daytime observations with $k_d > 0.3$. Additionally, a short description for each individual station is included below. 

\subsubsection{Bondville, IL (BON)}
The Bondville, Illinois SURFRAD station (BON) is located at $40.1^{\circ}$ N and $88.4^{\circ}$ W at an altitude of $\SI{213}{\meter}$. It was originally installed in April 1994 and uses a ventilated Eppley pyrgeometer to measure DLR. Close to Champaign, Illinois, the station has a Köppen climate classification of \textit{Dfa} (\textit{D} = continental climate, \textit{f} = no dry season, \textit{a} = hot summer). 

\subsubsection{Desert Rock, NV (DRA)}
The Desert Rock SURFRAD station (DRA) is located in the Mojave Desert at the Desert Rock operational radiosonde station at the Nevada Test Site, 65 miles outside of Las Vegas, Nevada. The precise coordinates are $36.6^{\circ}$ N and $116.0^{\circ}$ W, and the altitude of the station is $\SI{1007}{\meter}$. DRA also uses a ventilated Eppley pyrgeometer to measure DLR. With a Köppen climate classification of \textit{BWk} (\textit{B} = arid, \textit{W} = desert, \textit{k} = cold), this station is the most arid of the SURFRAD stations and experiences the least year-round cloud cover. 

\subsubsection{Goodwin Creek, MS (GWN)}
The Goodwin Creek SURFRAD station (GWN) outside Batesville, Mississippi is located at $34.3^{\circ}$ N and $89.9^{\circ}$ W with an altitude of $\SI{98}{\meter}$. This station also uses a ventilated Eppley pyrgeometer to measure DLR. Goodwin Creek is a watershed area in rural farmland in the southeast United States, belonging to the \textit{Cfa} (\textit{C} = temperate climate, \textit{f} = no dry season, \textit{a} = hot summer) classification of the Köppen climate scheme. On average, GWN is the most absolutely humid of the SURFRAD stations as well as the closest to sea level. 

\subsubsection{Fort Peck, MT (FPK)}
The Fort Peck, Montana SURFRAD station (FPK) is on the Fort Peck Tribes Reservation, outside of Polar, Montana. The station lies at $48.3^{\circ}$ N and $105.1^{\circ}$ W, with an altitude of $\SI{634}{\meter}$. The station was established in 1994 and uses a vented Eppley pyrgeometer for DLR measurements. Fort Peck has a Köppen climate classification of \textit{Cfa} (\textit{C} = continental climate, \textit{f} = no dry season, \textit{a} = hot summer). 

\subsubsection{Penn. State University, PA (PSU)}
The Penn. State University SURFRAD station (PSU) is located on the campus of Pennsylvania State University in State College, Pennsylvania at the university's agricultural research farm. Specifically, the coordinates of this location are $40.7^{\circ}$ N and $77.9^{\circ}$ W, and the station lies at an altitude of $\SI{376}{\meter}$. The farm sits in an Appalachian valley and experiences the Köppen climate designation of \textit{Dfb} (\textit{D} = continental climate, \textit{f} = no dry season, \textit{b} = warm summer).  

\subsubsection{Sioux Falls, SD (SXF)} 
The Sioux Falls SURFRAD station (SXF) is located at $43.7^{\circ}$ N and $96.6^{\circ}$ W with an altitude of $\SI{473}{\meter}$, which lies outside Sioux Falls, South Dakota. This station has a Köppen climate classification of \textit{Dfa}, which is the same as that of the Bondville, Illinois station. 

\subsubsection{Table Mountain, CO (TBL)}
The Table Mountain, Colorado SURFRAD instruments (TBL) are located at the Table Mountain Test Facility outside of Boulder, Colorado. The coordinates for this station are $40.1^{\circ}$ N and $105.2^{\circ}$ W, and the altitude is $\SI{1689}{\meter}$. These instruments were originally installed in July 1995. DLR is measured using a ventilated Eppley pyrgeometer. Like the station at Fort Peck, Montana, the Köppen climate classification for this station is \textit{BSk}. Located in the Rocky Mountains, this is the highest altitude SURFRAD station. 


\begin{figure}[htbp]% placement specifier
\centering %% For centre alignment of image.
\includegraphics[width=1\textwidth, keepaspectratio]{paper_pic/Surfard Location Map.png}
\caption{SURFRAD station locations. The seven atmospheric radiation monitoring stations are located across the continental United States, capturing data from a broad range of geographical settings and climates.} \label{fig:surfrad_map}
\end{figure}

\begin{table}[htbp]
\centering
\caption{Summary information of SURFRAD stations, including location, meteorology, and climate during 2010-2022.}
\label{tab:station-data}
\small 
\begin{tabular}{lccccccc}
\toprule
Station Name                   & BON   & DRA    & FPK   & GWN    & PSU   & SXF   & TBL   \\ 
\midrule
Latitude ($^\circ$)            & 40.1  & 36.6   & 48.3  & 34.2   & 40.7  & 43.7  & 40.1  \\
Longitude ($^\circ$)           & -88.4 & -116.0 & -105.1 & -89.9 & -77.9 & -96.6 & -105.2\\
Altitude (m)                   & 213   & 1007   & 634   & 98     & 376   & 473   & 1689  \\
Köppen classification          & \textit{Dfa} & \textit{BWk} & \textit{BSk} & \textit{Cfa} & \textit{Dfb} & \textit{Dfa} & \textit{BSk} \\ 
Average $T_a$ (℃)          & 17.1  & 24.7   & 15.2  & 23.2   & 15.0  & 12.9  & 15.6  \\
%25th Percentile of $T_a$ (℃) & 7.6   & 17.6   & 6.3   & 17.4   & 5.3   & 2.2   & 7.3   \\
%75th Percentile of $T_a$ (℃) & 27.1  & 32.5   & 25.5  & 30.5   & 24.5  & 24.4  & 24.9  \\
Average $P_w$ (hPa)        & 11.6  & 5.3    & 7.4   & 15.9   & 10.3  & 9.6   & 5.9   \\
%25th Percentile of $P_w$ (hPa) & 5.3   & 3.6    & 4.5   & 7.8    & 4.4   & 4.3   & 3.6   \\
%75th Percentile of $P_w$ (hPa) & 17.2  & 6.4    & 10.1  & 23.0   & 15.6  & 14.4  & 7.6   \\
Average DLR ($\SI{}{\watt \per \meter^2}$)       & 328.7 & 324.8  & 306.4 & 367.4  & 323.6 & 306.1 & 294.3 \\
%25th Percentile of DLR ($\SI{}{\watt \per \meter^2}$) & 272.9 & 289.0  & 265.3 & 316    & 275.2 & 251.0 & 256.1 \\
%75th Percentile of DLR ($\SI{}{\watt \per \meter^2}$) & 385.3 & 359.5  & 353.4 & 420.6  & 375.8 & 364.7 & 335.1 \\
Cloud frequency (\%)           & 33.0  & 13.2   & 33.6  & 28.7   & 42.5  & 32.0  & 27.0  \\
Data count               & 826,098 & 829,685 & 755,299 & 890,876 & 700,952 & 783,784 & 795,690 \\
\bottomrule
\end{tabular}
\end{table}

\subsection{Data Quality Control}

To ensure data quality and model accuracy, several preprocessing steps were undertaken before analysis and validation. These are listed below. 

\subsubsection{Radiation Data}

For quality control purposes, SURFRAD stations flag nonphysical data such as abnormally high or low temperatures and negative measured pressure and radiation values. All flagged data have been filtered from the dataset. Measurements have also been removed from the dataset if they exhibit negative DLR or $k_d$ values greater than 1. A lower bound is enforced on the GHI clearness index (data excluded if $k_{t,\mathrm{GHI}}<0.1$) to remove data points where the measured GHI is negligibly small compared to the clear-sky values.  

Several studies have formulated further limits for high-quality data. Inman et al. \cite{inman2016cloud} suggests data with high solar zenith angles ($\theta_z \geq 72.5^\circ$) should be excluded due to the low radiometer cosine response in this regime. These authors also show that while cloudy conditions can enhance GHI above clear-sky values, $k_t$ cannot exceed values of 1.5. Matsunobu and Coimbra \cite{matsunobu2024effective} found abnormalities in early morning measurements due to dew water having formed on the instrument, retaining data measured only after 8:00 local solar time. Long and Shi \cite{long2008automated} formulate several rare limits for shortwave measurements to be followed in real-world conditions, summarized here:  
\begin{equation}
    \mathrm{GHI} < 1.2 \,G_{on} \, (\cos\theta_z)^{1.2} + \SI{50}{\watt \per \meter^2},
\end{equation}

\begin{equation}
     \mathrm{DNI} < 0.95 \, G_{on} \, (\cos\theta_z)^{0.2} + \SI{10}{\watt \per \meter^2} ,
\end{equation}
where $G_{on}$ is the extraterrestrial irradiance on a surface normal to the sun's rays, adjusted for Earth-Sun distance, and $\theta_z$ is the solar zenith angle. These studies' recommendations have all been implemented in the SURFRAD data filtering for this study.

% Additionally, all observations are required to have positive downwelling longwave radiation ($\mathrm{DLR} > 0$), and the diffuse fraction must satisfy $k_d = \mathrm{DHI}/\mathrm{GHI} \leq 1$ to ensure physical consistency as $\mathrm{DHI} > \mathrm{GHI} $ violates the definition of diffuse fraction. A lower bound on the clearness index of $k_{t,\mathrm{GHI}} \geq 0.1$ is also imposed to exclude data points where the measured GHI is negligibly small relative to the clear-sky value, which can produce unreliable clearness index values. 

% The SURFRAD stations also pre-filter certain radiative data that cannot be sufficiently trusted, such as negative values for GHI, DNI, and DHI, and abnormally high or low temperature values (data are accepted only when $-80^\circ\mathrm{C} \leq T_a \leq 90^\circ\mathrm{C}$ and both the saturation vapor pressure $P_s$ and partial pressure of water vapor $P_w$ are non-negative).

% Additionally quality control flag filtering was performed as the SURFRAD dataset provided us with instrument-level quality control (QC) flags. The code requires all 7 QC flags (\texttt{qc\_direct\_n}, \texttt{qc\_dwsolar}, \texttt{qc\_diffuse}, \texttt{qc\_dwir}, \texttt{qc\_temp}, \texttt{qc\_rh}, \texttt{qc\_pressure}) to equal 0. The flag value of 0 indicates the highest data quality, and thus filtering ensures that only measurements meeting quality assurance standards are used in the analysis.

Additionally, NOAA has explicitly identified that the DRA site exhibited data anomalies beginning in 2017. Several studies that utilized this dataset have also reported issues during this period \citep{matsunobu2024effective,li2017determination}. In this study's analysis, it was observed that prediction errors increased significantly for DRA data collected after 2017 when compared to other sites and earlier years.
%   - DRA pre-2017  (n = 829,685): RMSE = 12.84 W/m^2, MAE = 8.52 W/m^2, Bias = +6.65 W/m^2
%   - DRA post-2017 (n = 757,837): RMSE = 41.46 W/m^2, MAE = 32.90 W/m^2, Bias = +32.66 W/m^2
%   - RMSE increase: +222.9%
%   - Other stations remained stable (max change: TBL +28.3%, BON +16.0%; FPK -2.2%, GWN +1.8%, PSU +0.1%, SXF -9.8%)
%Specifically, comparing clear-sky models residuals before and after 2017 at DRA reveals a pre-2017 RMSE of 12.84~W/m$^2$ versus a post-2017 RMSE of 41.46~W/m$^2$---a 223\% increase. The post-2017 mean bias of +32.66~W/m$^2$ indicates a systematic positive shift in measured DLR relative to clear-sky predictions. In contrast, the remaining six stations exhibited stable error metrics throughout 2010--2022 (see Fig.~\ref{fig:error_heatmaps}), confirming that this degradation is unique to DRA.
Therefore, all Desert Rock data collected after 2017 are excluded from the analysis. Further discussion surrounding station data reliability before and after 2017 is included in Appendix B. 

\subsubsection{Precipitation Data Filtration}
The effect of precipitation on the accuracy of radiation measurements remains unclear, with ongoing debate in the literature \citep{xu2014effect,wang2025rain}. Water droplets on sensor surfaces can interfere with accurate measurement of both shortwave and longwave radiation components. To mitigate potential sensor wetting artifacts, days with non-zero daily-average precipitation are excluded from the training dataset. The precipitation data used in this study is sourced from NASA's Prediction Of Worldwide Energy Resource (POWER) project, which provides daily global precipitation data. For each station, the daily average precipitation is computed, and only dates with a daily average of zero are retained for analysis.
%   - Date-level: 16,807 of 30,649 dates removed = 54.8% overall
%     Per station: BON 56.6%, DRA 27.7%, FPK 55.2%, GWN 57.4%, PSU 61.3%, SXF 54.4%, TBL 58.5%
%   - Observation-level: 12,262,678 of 22,101,313 observations removed = 55.5% overall

% This filter let to a significant reduction in the available dataset, with approximately 55\% of observation dates (55.5\% logged data points) removed across all stations. The removal rate varied by station, with the arid DRA site having the lowest removal rate of 27.7\% while the more humid PSU site had the highest removal rate of 61.3\%. While the following step reduces the fitting dataset, it is necessary to ensure that the model is trained on data that are not affected by precipitation-induced measurement errors, which could otherwise bias the model and reduce its generalizability.

\subsubsection{Post-Calculation Data Filtering}

After computing the derived quantities $\varepsilon_{\mathrm{sky}}$, $\varepsilon_{\mathrm{clear\,sky}}$, and $\gamma$, additional data cleaning is performed. Any observations leading to non-finite values (e.g. due to divison by zero in the $\gamma$ calculation when $\varepsilon_{\mathrm{clear\,sky}} \approx 1$) or values of $\gamma > 1$ are removed. % cloud fraction values are filtered out by excluding observations with $\gamma > 1$. This is because values of $\gamma$ exceeding unity would imply an effective sky emissivity greater than that of a perfect blackbody, which is not physically meaningful in this context. 

\subsubsection{Balancing Sample Sizes for Model Training}

Since the dataset used in this study includes observations from different radiation stations, the amount of data retained after applying the aforementioned filtering criteria varies among stations. To minimize potential station biases in the model fitting process and to ensure that the resulting model is not disproportionately influenced by stations with larger datasets, we adopt a uniform data volume approach. To do this, the station with the smallest remaining dataset after filtering is selected as a reference, and all other stations with greater data volumes are randomly downsampled so that each station's dataset size matches this reference. This approach ensures that each station contributes equally to the model training process, thereby enhancing the generalizability and robustness of the fitted model. The initial data counts, after radiation, precipitation, and post-calculation filtering, are shown in Table~\ref{tab:station-data}. After downsampling, each station retains 700,952 data points.  

% \begin{table}[!htbp]
% \centering
% \caption{Number of SURFRAD data points after filtering and final count used for model fitting }
% \label{tab:data-counts}
% \small 
% \begin{tabular}{lccccccc}
% \toprule
% station & after data quality filtering & after station bias filtering \\
% \midrule
% BON     & 826,098                       & 700,952                       \\
% DRA     & 829,685                       & 700,952                       \\
% FPK     & 755,299                       & 700,952                       \\
% GWN     & 890,876                       & 700,952                       \\
% PSU     & 700,952                       & 700,952                       \\
% SXF     & 783,784                       & 700,952                       \\
% TBL     & 795,690                       & 700,952                      \\
% \bottomrule
% \end{tabular}
% \end{table}

\subsection{Error Metrics}
Based on the dataset with the aforementioned filtering conditions applied, the models listed in Eq.s~\eqref{model_1}-\eqref{model_3} are fit and validated with respect to several error metrics. The error metrics used in this study are the mean bias error (MBE), the root mean square error (RMSE), the reduced mean bias error (rMBE), the reduced root mean square error (rRMSE), and the coefficient of determination ($R^2$), and their formulae are listed below.

\begin{enumerate}
    \item Mean bias error (MBE):
    \begin{equation}
        \mathrm{MBE} = \frac{1}{n} \sum_{i=1}^{n} \left( y_{\mathrm{observed}}^i - y_{\mathrm{predicted}}^i \right),
    \end{equation}
    where $y_{\mathrm{observed}}$ is the observed value, $y_{\mathrm{predicted}}$ is the predicted value, and $n$ is the total number of data points.

    \item Root mean squared error (RMSE):
    \begin{equation}  
        \mathrm{RMSE} = \sqrt{\frac{1}{n} \sum_{i=1}^{n} \left( y_{\mathrm{observed}}^i - y_{\mathrm{predicted}}^i \right)^2}.
    \end{equation}
    This metric is used to measure the deviation between the predicted and observed values. A smaller value indicates a better fit of the model. 

    \item Relative mean bias error (rMBE):
    \begin{equation}
        \mathrm{rMBE} = \frac{1}{n} \sum_{i=1}^{n} \frac{y_{\mathrm{observed}}^i - y_{\mathrm{predicted}}^i}{y_{\mathrm{observed}}^i}.
    \end{equation}
    The rMBE considers the relative deviation of the data, providing a better comparison of model performance across different scales.

    \item Relative root mean squared error (rRMSE):
    \begin{equation}
        \mathrm{rRMSE} = \sqrt{\frac{1}{n} \sum_{i=1}^{n} \left( \frac{y_{\mathrm{observed}}^i - y_{\mathrm{predicted}}^i}{y_{\mathrm{observed}}^i} \right)^2}.
    \end{equation}
    The rRMSE helps assess the model's fitting performance for observed values of different magnitudes.

    \item Coefficient of determination ($R^2$)
    \begin{equation}
        R^2 = 1 - \frac{\sum_{i=1}^n \left(y_{\mathrm{observed}}^i - y_{\mathrm{predicted}}^i \right)^2}{\sum_{i=1}^n \left(y_\mathrm{observed}^i - \bar{y} \right)^2},
    \end{equation}
    where $\bar{y}$ is the mean of the observed data. The $R^2$ value measures how much data variance is attributable to the model, with values closer to $1$ indicating a more robust model. 

\end{enumerate}

\section{Results and Discussion}

\subsection{Determination of model parameters}

The coefficients for each proposed model in Eq.s~\eqref{model_1}-\eqref{model_3} were found to minimize the root mean squared error (RMSE) between the observed cloud fraction factor $\gamma$ from the SURFRAD data and the predicted cloud fraction factor for each model. Results for the RMSE for a range of possible coefficient values are shown in the contour plots in Fig.~\ref{fig:rmse_contour}. Model 1, corresponding to the diffuse fraction $k_d$, achieves the lowest total RMSE of 0.0296 for coefficients of $c_1=0.585$ and $c_2=1.748$, indicating the highest accuracy of all three models. Model 2 and Model 3, corresponding to the clear-sky $k_t$ indices, both demonstrate lower RMSE values and thus lower predictive performance. 

\begin{figure}[htbp]% placement specifier
\centering%% For centre alignment of image.
\includegraphics[width=1\textwidth, keepaspectratio]{paper_pic/RMSE_contour.png}
\caption{Contour plots showing the root mean squared error (RMSE) for each proposed model with respect to fitting parameters $c_1$ and $c_2$. The fitting parameters for each model are chosen to minimize RMSE, with the minimum point indicated on each plot by the black triangle and its coordinates labeled in white. (a) Model 1, function of diffuse fraction, $k_d$. This model achieves the lowest overall RMSE. (b) Model 2, function of clearness index based on GHI, $k_{t,\mathrm{GHI}}$. (c) Model 3, function of clearness index based on DNI, $k_{t,\mathrm{DNI}}.$
} \label{fig:rmse_contour}
\end{figure}

A comparison of the three models' predictive performance is shown in the box plots in Fig.~\ref{fig:model_boxplots}, which display the distribution of prediction errors for both emissivity and DLR across the three models. Model 1 achieves the lowest median bias, tightest error distribution, and lowest interquartile range. Because Model 1 showcases the greatest accuracy and predictive value, this is chosen as the proposed overall radiation model: 

\begin{equation} \label{final_model}
    \gamma = 0.585 \, k_d^{1.748} \, .
\end{equation}

%We can further compare the performance of the three models by analyzig the box plots in Fig.~\ref{fig:model_boxplots}, which show the distribution of prediction errors for both emissivity and DLR across the three models. Model~1, which is based on the diffuse fraction $k_d$, exhibits the tightest error distribution, with the smallest median bias and interquartile range (IQR). In contrast, Models~2 and 3, which are based on the clearness indices $k_t$, show wider spreads in their error distributions and larger systematic errors. This analysis further supports the conclusion that the diffuse fraction is a more reliable parameterization of cloud effects on longwave emissivity compared to the clearness indices.

\begin{figure}[htbp]% placement specifier
\centering%% For centre alignment of image.
\includegraphics[width=0.85\textwidth, keepaspectratio]{paper_pic/model_comparison_boxplots.png}
\caption{Boxplots comparing prediction errors across the three models for (a) all-sky emissivity and (b) downwelling longwave radiation. Each box displays the median prediction error and the interquartile range, with whiskers representing the $1.5\times$ interquartile range and additional points included as outliers. Model 1 displays the lowest median error and tightest distribution among the three models for both emissivity and DLR predictions. }
\label{fig:model_boxplots}
\end{figure}

% We can attribute various physical reasons to the superior performance of the diffuse fraction $k_d$ over the clearness indices $k_t$ for parameterizing cloud effects on downwelling longwave radiation (DLR). First, the clearness index ($k_t$) conflates cloud attenuation with aerosol effects. Both clouds and aerosols leads to the reduction of $k_t$, but aerosols do not enhance DLR the way an optically thick cloud would. This ambiguity would lead to added noise in the $k_t$--$\gamma$ relationship, as the model can't distinguish between low $k_t$ values caused by clouds versus those caused by aerosols. However, the diffuse fraction ($k_d$) directly measures the proportion of incoming solar radiation that has been scattered by clouds. A high $k_d$ value indicates a diffuse-dominated sky, which corresponds to cloudy conditions and thus a high $\gamma$. This provides a more physically direct link to cloud radiative properties, making $k_d$ a more reliable parameter for modeling cloud effects on longwave emissivity.

\subsection{Station-specific trends}
Equation~\eqref{final_model} reflects the overall radiative model which was fit to data from all seven SURFRAD stations. However, to more specifically analyze the accuracy of Eq.~\eqref{final_model} applied to different stations and climate patterns, parameters for the general $k_d$ model from Eq.~\eqref{model_1} were also fit with respect to each individual SURFRAD station. Results for each station, as well as the overall fit, are shown in Fig.~\ref{fig:kd_validation}. 

\begin{figure}[!htbp]% placement specifier
%% Use \includegraphics command to insert graphic files. Place graphics files in 
%% working directory.
\centering%% For centre alignment of image.
\includegraphics[width=0.85\textwidth, keepaspectratio]{paper_pic/all_density_plots_labeled.png}
%% Use \caption command for figure caption and label.
  \caption{Data density plots illustrating the relationship between diffuse fraction $ k_{d}$ and cloud factor $ \gamma$ in the SURFRAD data, overlaid with best-fit power-law models. (a) Data density and overall power-law fit for all stations. The overall model, given by Eq.~\eqref{final_model}, is displayed in orange on all subplots for comparison. (b-h) Data density and power-law fits for individual stations Bondville (BON), Desert Rock (DRA), Fort Peck (FPK), Goodwin Creek (GWN), Pennsylvania State University (PSU), Sioux Falls (SXF), and Table Mountain (TBL), respectively. Individual power-law fits are chosen to minimize RMSE for each station's data, with the equations for the fits displayed in the subplot legends and the station models overlaid in cyan. Each individual station model, excepting DRA, shows good agreement with the overall model.  
  }
    \label{fig:kd_validation}
\end{figure}

Each subplot in Fig.~\ref{fig:kd_validation} visualizes the site-specific ${k_{d}}$–$\gamma$ relationship, highlighting both the general trend and unique station-level variations due to local climatic and atmospheric conditions. The station-specific fits for BON, FPK, GWN, PSU, SXF, and TBL match well with the general trend, whereas the fit for DRA deviates significantly. 

% Table~\ref{tab:station-climate-fit} summarizes the key climatological properties and model-fit error metrics for each station. The stations span a wide range of atmospheric conditions: GWN, the lowest-altitude site (\SI{98}{\meter}) with the highest mean water vapor pressure ($\overline{P_w}=14.64$~hPa), exhibits the highest clear-sky emissivity ($\overline{\varepsilon}_{\mathrm{clear\,sky}}=0.788$) and the smallest emissivity gap ($1-\overline{\varepsilon}_{\mathrm{clear\,sky}}=0.212$). Consequently, cloud-induced enhancements in $\gamma$ are moderate at GWN, and its station-specific fit aligns closely with the global power-law, yielding the lowest DLR rRMSE of 2.84\%. In contrast, TBL---the highest-altitude station (\SI{1689}{\meter}) with low precipitable water ($\overline{P_w}=5.80$~hPa)---has the lowest clear-sky emissivity ($\overline{\varepsilon}_{\mathrm{clear\,sky}}=0.693$) and the largest emissivity gap (0.307). This wider gap amplifies the relative cloud enhancement and produces a slight offset from the global fit at intermediate $k_d$ values, resulting in the highest DLR rRMSE of 4.62\%.

% The four mid-latitude continental stations (BON, FPK, PSU, SXF) share similar temperate climates with mean temperatures between 14.6 and \SI{16.9}{\celsius} and mean $P_w$ between 7.47 and 11.42~hPa. Their broad seasonal cloud variability---reflected in cloud frequencies ranging from 32.0\% to 42.5\%---provides a wide range of $k_d$ values for fitting, and all four station-specific fits closely track the global power-law. Among these, PSU exhibits the highest cloud frequency (42.5\%) and the largest DLR rRMSE (4.08\%), which may relate to localized Appalachian valley cloudiness effects.

% \begin{table}[!htbp]
% \centering
% \caption{Station climatological properties and model-fit error metrics (training period, 2010--2022). Cloud frequency is defined as the percentage of daytime observations with $k_d > 0.3$. Model errors are computed using the global power-law fit (Eq.~\ref{final_model}).}
% \label{tab:station-climate-fit}
% \small
% \begin{tabular}{lccccccc}
% \toprule
%  & \textbf{BON} & \textbf{DRA} & \textbf{FPK} & \textbf{GWN} & \textbf{PSU} & \textbf{SXF} & \textbf{TBL} \\
% \midrule
% Altitude (m)                          & 213   & 1007  & 634   & 98    & 376   & 473   & 1689  \\
% Mean $T_a$ (\si{\celsius})            & 16.9  & 23.8  & 14.9  & 22.0  & 15.1  & 14.6  & 15.1  \\
% Mean $P_w$ (hPa)                      & 11.42 & 4.97  & 7.47  & 14.64 & 9.89  & 10.51 & 5.80  \\
% $\overline{\varepsilon}_{\mathrm{clear\,sky}}$ & 0.762 & 0.696 & 0.725 & 0.788 & 0.748 & 0.750 & 0.693 \\
% $1-\overline{\varepsilon}_{\mathrm{clear\,sky}}$ & 0.238 & 0.304 & 0.275 & 0.212 & 0.253 & 0.250 & 0.307 \\
% Cloud freq.\,(\%)                     & 33.0  & 13.2  & 33.6  & 28.7  & 42.5  & 32.0  & 27.0  \\
% \midrule
% $\gamma$ RMSE                         & 0.114 & 0.076 & 0.112 & 0.121 & 0.135 & 0.123 & 0.113 \\
% DLR RMSE (W/m$^2$)                    & 10.22 & 9.88  & 11.39 & 10.04 & 12.66 & 11.06 & 13.09 \\
% DLR rRMSE (\%)                        & 3.19  & 3.13  & 3.77  & 2.84  & 4.08  & 3.60  & 4.62  \\
% \bottomrule
% \end{tabular}
% \end{table}

Compared to the other stations, DRA experiences little annual cloud cover due to its arid climate (from Table~\ref{tab:station-data}, DRA is the only station in the "arid" Köppen climate classification group \textit{B} and also has the lowest value of average $P_w$). This is reflected by the lack of data at higher values of $k_d$ and $\gamma$ as seen in Fig.~\ref{fig:kd_validation}. As such, the station-specific model for DRA is fit to data with $k_d$ and $\gamma$ values close to zero, which do not capture the all-sky conditions in the overall all-station model. Because of the disparity between the general fit and the station-specific DRA fit reflecting DRA's uniquely arid conditions, the DRA fit is excluded from further validation analysis. Further comparisons and justifications for the exclusion of the DRA fit are included in Appendix A.  

\subsection{Uncertainty and Prediction Intervals}
Uncertainty analysis is also conducted to quantify the confidence in the model predictions. A bootstrap resampling approach is employed, in which the training dataset is resampled with replacement 1000 times to generate distributions of the fitted parameters $c_1$ and $c_2$. The standard deviation of the residuals from the original fit is used to compute prediction intervals of the cloud fraction factor $\gamma$. The resulting 68\% and 95\% prediction bands are visualized in Fig.~\ref{fig:uncertainty}, with the best-fit curve from Eq.~\eqref{final_model} shown in dark blue and the underlying data density represented in grayscale. The residuals are found to be approximately Gaussian, and the empirical coverage rates closely matched theoretical expectations, with approximately 68\% of observations falling within $\pm 1\sigma$ and 95\% within $\pm 2\sigma$. Thus, the analysis confirms that the model uncertainty is well-characterized and the prediction intervals can be reliably used for applications such as climate modeling and solar forecasting.

\begin{figure}[!htbp]
\centering
\includegraphics[width=0.85\textwidth, keepaspectratio]{paper_pic/uncertainty_prediction_bands.png}
\caption{Plot of the model prediction uncertainty. The proposed model is shown as the solid line, with the dark and light shaded regions representing the $\pm 1\sigma$ and $\pm 2\sigma$ prediction confidence intervals, respectively.}
\label{fig:uncertainty}
\end{figure}

\subsection{Validation using 2023-2024 observations}

To evaluate the robustness and generalizability of the proposed model, it is validated using more recent SURFRAD datasets from 2023 and 2024. The statistical evaluation results are summarized in Table~\ref{tab:validation}. The table reports the mean bias error (MBE), root mean square error (RMSE), relative MBE (rMBE), relative RMSE (rRMSE), and coefficient of determination ($R^2$) for each station.

\begin{table}[htbp]
\centering
\caption{Validation results using 2023\textendash 2024 SURFRAD data (excluding DRA).}
\label{tab:validation}
\begin{tabular}{lcccccc}
\hline
\textbf{Station} & \textbf{MBE} & \textbf{RMSE} & \textbf{rMBE} & \textbf{rRMSE} & $\mathbf{R^2}$ & \textbf{p-value} \\
\hline
BON  & 0.0031  & 0.0349  & 0.376\% & 4.262\% & 0.8221  & $< 0.01$ \\
FPK  & 0.0075  & 0.0369  & 0.948\% & 4.657\% & 0.7687  & $< 0.01$ \\
GWN  & 0.0002  & 0.0297  & 0.018\% & 3.532\% & 0.8752  & $< 0.01$ \\
PSU  & 0.0066  & 0.0433  & 0.789\% & 5.179\% & 0.8142  & $< 0.01$ \\
SXF  & 0.0043  & 0.0356  & 0.534\% & 4.405\% & 0.8296  & $< 0.01$ \\
TBL  & 0.0065  & 0.0389  & 0.862\% & 5.172\% & 0.7572  & $< 0.01$ \\
\textbf{All stations} & \textbf{\,0.0047} & \textbf{\,0.0369} & \textbf{\,0.581\%} & \textbf{\,4.563\%} & \textbf{\,0.8302} & \textbf{$< 0.01$} \\
\hline
\end{tabular}
\end{table}

Overall, the results demonstrate that the model maintains high predictive accuracy and consistency across most stations, with low bias and strong correlations between observed and predicted values. The average MBE across all stations is approximately 0.0047, indicating a minimal systematic deviation. The rMBE remains below 1\% for all stations, suggesting that the model is not prone to significant over- or underestimation.

The RMSE and rRMSE values range from 0.0297 to 0.0433 and 3.53\% to 5.18\%, respectively, reflecting acceptable levels of dispersion around the predicted values. The highest rRMSE values occur at PSU and TBL, both reaching 5.18\%, which may be attributed to localized climatic anomalies (e.g. the location of PSU in an Appalachian valley and the high altitude of TBL). On the other hand, the lowest error is observed at GWN, with an RMSE of 0.0297 and rRMSE of 3.53\%, demonstrating strong model performance at this site.

The $R^2$ values are consistently high, exceeding 0.75 at all stations, with a maximum of 0.875 at GWN. This confirms a strong correlation between the model predictions and actual observations, reinforcing the model's applicability across diverse geographical conditions. Furthermore, all p-values are below 0.01, indicating that the correlations between predicted and observed values are statistically highly significant across all stations. This reinforces the robustness and reliability of the model's predictive capabilities.

Collectively, these results indicate that the model maintains robust predictive performance and stability when applied to new datasets beyond the initial training period.

\subsubsection{Validation using estimated $k_d$ derived from $ k_t$}

Without measurements of DHI, the diffuse fraction $k_d$ cannot be computed and thus the proposed radiative model cannot be directly applied. This leads to challenges, since many radiation monitoring stations outside the SURFRAD network are not equipped to measure diffuse solar radiation components, opting instead to measure only GHI due to the lower instrumentation cost and simplicity of installation \cite{nollas2023quality,marion2015model}. However, several empirical models are available to estimate the diffuse fraction $k_d$ from the clearness index $k_{t,\mathrm{GHI}}$, which can be computed from GHI measurements only. All-sky emissivity can then be indirectly found using the estimated $k_d$ values from these models and the radiative model proposed in this study.  

Below, we list three empirical clearness index models from \cite{erbs1982estimation}, \cite{orgill1977correlation}, and \cite{reindl1990diffuse}, and validate the results obtained from the estimated $k_d$ values using SURFRAD data. The models are as follows:

\begin{enumerate}

\item \textbf{Orgill–Hollands (1977) Model:}
\begin{equation}
k_d = \begin{cases}
1.0 - 0.249 k_t & k_t \leq 0.35 \\
1.557 - 1.84 k_t & 0.35 < k_t \leq 0.75 \\
0.177 & k_t > 0.75
\end{cases}
\end{equation}

\item \textbf{Erbs et al. (1982) Model:}
\begin{equation}
k_d = \begin{cases}
1.0 - 0.09 k_t & k_t \leq 0.22 \\
0.9511 - 0.1604 k_t + \\
4.388 k_t^2 - 16.638 k_t^3 + 12.336 k_t^4 & 0.22 < k_t \leq 0.8 \\
0.165 & k_t > 0.8
\end{cases}
\end{equation}

\item \textbf{Reindl et al. (1990) Model:}
\begin{equation}
k_d = \begin{cases}
1.45 - 1.67 k_t, \quad 0 \leq k_t \leq 0.3 \\
1.02 - 1.02 k_t, \quad 0.3 < k_t \leq 0.78 \\
0.147, \quad k_t > 0.78
\end{cases}
\end{equation}

\end{enumerate}

To facilitate the comparison of measured SURFRAD data with the predicted all-sky emissivity values calculated from the empirical models above, the all-sky emissivity is computed as a function of the clearness index $k_t$ and the dimensionless precipitable water $p_w$. This visualization is shown in Fig.~\ref{fig:kd_validation_a} for each of the three empirical models. Visual inspection of Fig.~\ref{fig:kd_validation_a} demonstrates agreement between the measured SURFRAD data, shown as gray points, and the predicted emissivity, shown as the graded surface, without noticeable systematic discrepancies. %Further details are included in Fig.~\ref{fig:kd_validation_1} and Fig.~\ref{fig:kd_validation_2}, which are the projections of the 3D emissivity surfaces onto the $\varepsilon_\mathrm{sky}$-$k_t$ and $\varepsilon_\mathrm{sky}$-$p_w$ planes, respectively.

\begin{figure}[!htbp]
\centering
\includegraphics[width=1\textwidth, keepaspectratio]{paper_pic/Combined_3D_e_sky_Surfaces.png}
  \caption{3D surfaces of all-sky emissivity calculated with the proposed $k_d$-$\gamma$ model assuming $k_d$ data is unavailable. Without $k_d$ data, emissivity may still be estimated using empirical models which relate $k_d$ to the more widely available clearness index, $k_t$. SURFRAD data is shown as the gray points, indicating reasonable agreement with emissivity surfaces. Models used here are: (a) Orgill-Hollands (1977); (b) Erbs et al. (1982); (c) Reindl et al. (1990a).}
    \label{fig:kd_validation_a}
\end{figure}

% \begin{figure}[!htbp]
% \centering
% \includegraphics[width=1\textwidth, keepaspectratio]{paper_pic/Combined_kt_e_sky_fixedCB.png}
%   \caption{Projection of 3D all-sky emissivity surfaces onto the $\varepsilon_\mathrm{sky}$-$k_t$ plane for the Orgill–Hollands (1977), Erbs et al.\ (1982), and Reindl et al.\ (1990a) models.}
%     \label{fig:kd_validation_1}
% \end{figure}

% \begin{figure}[!htbp]
% \centering
% \includegraphics[width=1\textwidth, keepaspectratio]{paper_pic/Combined_pw_e_sky_fixedCB.png}
%   \caption{Projection of 3D all-sky emissivity surfaces onto the $\varepsilon_\mathrm{sky}$-$p_w$ plane for the Orgill–Hollands (1977), Erbs et al.\ (1982), and Reindl et al.\ (1990a) models.}
%     \label{fig:kd_validation_2}
% \end{figure}

For a more quantitative validation, error metrics are calculated between the measured longwave emissivity values and the emissivity values predicted using GHI measurements, the empirical $k_t$ models, and the radiative model proposed in this study. Results are listed in Table~\ref{tab:validation_statistics} across different stations. As shown, the relative mean bias error (rMBE) ranges from 2.7\% to 6.1\% for all models and stations, with the \cite{erbs1982estimation} model showing slightly lower bias at most stations and the \cite{reindl1990diffuse} model having slightly higher rMBE values.

All models show rRMSE values between 5.4\% and 10.0\%. The highest errors appear at station FPK for all three models, which may be due to lower accuracy in $k_d$ estimation. In contrast, station GWN shows the lowest rRMSE values, which is consistent with its performance in direct validation. This suggests that the models remain reliable when applied using estimated inputs. 

While these error metrics show slightly lower performance than when using directly calculated $k_d$ values, the model retains reasonable predictive power. This highlights the practical utility of the model in observationally constrained environments.

\begin{table}[!htbp]
\centering
\caption{Statistical validation of sky emissivity models across different stations}
\label{tab:validation_statistics}
\renewcommand{\arraystretch}{1.2}
\begin{tabular}{llcccc}
\toprule
\textbf{Station} & \textbf{Model} & \textbf{MBE} & \textbf{RMSE} & \textbf{rMBE} & \textbf{rRMSE} \\
\midrule
\multirow{3}{*}{BON} & Erbs et al. (1982)          & 0.0350 & 0.0588 & 0.0423 & 0.0711 \\
                     & Orgill–Hollands (1977)      & 0.0350 & 0.0593 & 0.0423 & 0.0716 \\
                     & Reindl et al. (1990a)       & 0.0364 & 0.0599 & 0.0440 & 0.0724 \\
\midrule
\multirow{3}{*}{FPK} & Erbs et al. (1982)          & 0.0468 & 0.0791 & 0.0581 & 0.0982 \\
                     & Orgill–Hollands (1977)      & 0.0468 & 0.0792 & 0.0581 & 0.0984 \\
                     & Reindl et al. (1990a)       & 0.0486 & 0.0802 & 0.0603 & 0.0996 \\
\midrule
\multirow{3}{*}{GWN} & Erbs et al. (1982)          & 0.0225 & 0.0460 & 0.0265 & 0.0542 \\
                     & Orgill–Hollands (1977)      & 0.0228 & 0.0469 & 0.0268 & 0.0553 \\
                     & Reindl et al. (1990a)       & 0.0239 & 0.0471 & 0.0282 & 0.0555 \\
\midrule
\multirow{3}{*}{PSU} & Erbs et al. (1982)          & 0.0371 & 0.0621 & 0.0442 & 0.0739 \\
                     & Orgill–Hollands (1977)      & 0.0380 & 0.0636 & 0.0452 & 0.0757 \\
                     & Reindl et al. (1990a)       & 0.0391 & 0.0638 & 0.0465 & 0.0760 \\
\midrule
\multirow{3}{*}{SXF} & Erbs et al. (1982)          & 0.0371 & 0.0659 & 0.0455 & 0.0807 \\
                     & Orgill–Hollands (1977)      & 0.0372 & 0.0664 & 0.0455 & 0.0814 \\
                     & Reindl et al. (1990a)       & 0.0387 & 0.0671 & 0.0474 & 0.0823 \\
\midrule
\multirow{3}{*}{TBL} & Erbs et al. (1982)          & 0.0361 & 0.0611 & 0.0470 & 0.0795 \\
                     & Orgill–Hollands (1977)      & 0.0362 & 0.0615 & 0.0472 & 0.0800 \\
                     & Reindl et al. (1990a)       & 0.0383 & 0.0623 & 0.0498 & 0.0810 \\
\bottomrule
\end{tabular}
\end{table}

% \subsection{Operational Implementation Framework}

% For estimating all-sky downwelling longwave radiation (DLR) in an operational context, the proposed model can be implmented using the procedure outlined in Algorithm~\ref{alg:dlr_estimation}. This algorithm provides a step-by-step framework for users to estimate DLR using commonly available meteorological and radiation measurements, such as air temperature, relative humidity, global horizontal irradiance (GHI), and either diffuse horizontal irradiance (DHI) or clear-sky GHI. By following this procedure, one can compute the necessary parameters to apply the discussed radiative model and obtain accurate estimates of all-sky DLR for various applications in climate research, weather forecasting, and energy modeling.

% \begin{algorithm}[!htbp]
% \caption{All-Sky Downwelling Longwave Radiation (DLR) Estimation}
% \label{alg:dlr_estimation}
% \begin{algorithmic}[1]
% \Require Air temperature $T_a$ [K], Relative humidity $\phi$ [0--1], Site altitude $z$ [m]
% \Require Global Horizontal Irradiance (GHI) [W/m$^2$]
% \Require Diffuse Horizontal Irradiance (DHI) [W/m$^2$] \textbf{or} Clear-sky GHI [W/m$^2$]
% \Ensure Downwelling Longwave Radiation DLR [W/m$^2$]
% \Statex
% \State \textbf{Constants:}
% \State \hspace{1em} $\sigma \gets 5.67 \times 10^{-8}$ W/(m$^2\cdot$K$^4$) \Comment{Stefan--Boltzmann constant}
% \State \hspace{1em} $P_0 \gets 101325$ Pa \Comment{Standard atmospheric pressure}
% \Statex
% \State \textbf{Step 1: Compute Saturation Vapor Pressure}
% \State \hspace{1em} $P_s \gets 610.94 \cdot \exp\left(\dfrac{17.625 \cdot (T_a - 273.15)}{T_a - 30.11}\right)$ \Comment{[Pa]}
% \Statex
% \State \textbf{Step 2: Compute Partial Pressure of Water Vapor}
% \State \hspace{1em} $P_w \gets P_s \cdot \phi$ \Comment{[Pa]}
% \State \hspace{1em} $p_w \gets P_w / P_0$ \Comment{Dimensionless}
% \Statex
% \State \textbf{Step 3: Compute Clear-Sky Emissivity}
% \State \hspace{1em} $\varepsilon_{\text{clear}} \gets 0.600 + 1.652 \cdot \sqrt{p_w} + 0.150 \cdot \left[\exp\left(-\dfrac{z}{8500}\right) - 1\right]$
% \Statex
% \State \textbf{Step 4: Determine Diffuse Fraction $k_d$}
% \If{DHI is measured}
%     \State \hspace{1em} $k_d \gets \text{DHI} / \text{GHI}$ \Comment{Direct calculation}
% \Else
%     \State \hspace{1em} Compute clearness index: $k_t \gets \text{GHI}_{\text{measured}} / \text{GHI}_{\text{clear-sky}}$
%     \State \hspace{1em} Estimate $k_d$ using Erbs et al.\ (1982):
%     \If{$k_t \leq 0.22$}
%         \State \hspace{2em} $k_d \gets 1.0 - 0.09 \cdot k_t$
%     \ElsIf{$k_t \leq 0.80$}
%         \State \hspace{2em} $k_d \gets 0.9511 - 0.1604 \cdot k_t + 4.388 \cdot k_t^2 - 16.638 \cdot k_t^3 + 12.336 \cdot k_t^4$
%     \Else
%         \State \hspace{2em} $k_d \gets 0.165$
%     \EndIf
% \EndIf
% \Statex
% \State \textbf{Step 5: Compute Cloud Fraction Factor}
% \State \hspace{1em} $\gamma \gets 0.585 \cdot k_d^{1.748}$
% \Statex
% \State \textbf{Step 6: Compute All-Sky Emissivity}
% \State \hspace{1em} $\varepsilon_{\text{sky}} \gets \varepsilon_{\text{clear}} + \gamma \cdot (1 - \varepsilon_{\text{clear}})$
% \Statex
% \State \textbf{Step 7: Compute Downwelling Longwave Radiation}
% \State \hspace{1em} $\text{DLR} \gets \varepsilon_{\text{sky}} \cdot \sigma \cdot T_a^4$ \Comment{[W/m$^2$]}
% \Statex
% \State \Return DLR
% \end{algorithmic}
% \end{algorithm}


% Finally, Table~\ref{tab:operational_guidance} provides operational guidance for selecting the appropriate DLR estimation approach based on the availability of measurements from a given station. The table summarizes the recommended method for each scenario, along with expected error metrics relative to measured DLR values. This guidance can assist users in choosing the most suitable estimation technique for their specific observational constraints and accuracy requirements.

% \begin{table}[!htbp]
% \centering
% \caption{Operational guidance for downwelling longwave radiation (DLR) estimation based on available measurements; expected errors relative to measured DLR.}
% \label{tab:operational_guidance}
% \small
% \begin{tabular}{p{3.5cm}p{2.5cm}p{4.8cm}cc}
% \toprule
% \textbf{Scenario} & \textbf{Available Data} & \textbf{Recommended Approach} & \textbf{rRMSE} & \textbf{$R^2$} \\
% \midrule
% A: Full radiation & GHI, DHI, $T_a$, $\phi$, $z$ & Direct $k_d = \text{DHI}/\text{GHI}$, Eq.~\eqref{final_model} & 4.6\% & 0.83 \\
% \addlinespace
% B: GHI only & GHI, $T_a$, $\phi$, $z$ & Estimate $k_d$ via Erbs et al.\ (1982), Eq.~\eqref{final_model} & 5.4--10\% & 0.70--0.80 \\
% \addlinespace
% C: Meteorology only & $T_a$, $\phi$, $z$ & Clear-sky model only (Eq.~\ref{e_clear sky}), no cloud correction & 15--20\% & -- \\
% \bottomrule
% \end{tabular}
% \end{table}

% \subsection{Sensitivity Analysis}

% To determine the sensitivity of the proposed model to its parameters ($c_1$ and $c_2$ in Eq.~\eqref{final_model}), a perturbation analysis was conducted. Each parameter was independently varied by $\pm 5\%$, $\pm 10\%$, and $\pm 15\%$ from their optimal values, and the resulting RMSE of the predicted cloud fraction factor $\gamma$ was evaluated against the full training dataset. The results of this analysis are visualized in Fig.~\ref{fig:sensitivity}. Simultaneous, positive perturbation in $c_1$ and negative perturbation in $c_2$ (opposite directions), led to relatively higher increase in RMSE values as shown in the lower left corners of the heat map.

% \begin{figure}[!htbp]
% \centering
% \includegraphics[width=0.75\textwidth, keepaspectratio]{paper_pic/sensitivity_analysis_heatmap.png}
% \caption{Heatmap of RMSE of cloud fraction factor $\gamma$ predictions as parameters $c_1$ and $c_2$ in the model are independently perturbed by $\pm 5\%$, $\pm 10\%$, and $\pm 15\%$ from their optimal values. The minimum RMSE occurs at the unperturbed center, indicating the sensitivity of the model to parameter variations.}
% \label{fig:sensitivity}
% \end{figure}

\section{Conclusion}

In this study, a novel parameterization model is proposed for predicting sky emissivity and downwelling longwave radiation under all-sky conditions by leveraging shortwave radiation measurements. The best performing model uses a power law function to relate the cloud fraction factor $\gamma$ to the diffuse fraction $k_d$. By taking data from seven high-quality data monitoring stations spaced across various climates and locations, the model accurately and robustly predicts the influence of clouds on the thermal energy balance at the surface.

Comprehensive validations were carried out with multiple strategies. Direct validation with recent observations yielded low rRMSE values (3.53\%-5.18\%), with station GWN exhibiting the smallest error. This confirms the model's consistency and robustness when applied to newly acquired data. Indirect validation was also conducted using values of $k_d$ estimated from widely used empirical models. Despite the reduced accuracy associated with estimated $k_d$, the model maintained acceptable performance, reinforcing its practical value when complete radiation data are unavailable.

Overall, these results highlight that even with limited input variables, the model maintains reliable performance across time and space. GWN consistently yielded the lowest errors in both validation modes, which may be attributed to its regional meteorological conditions. Importantly, this work demonstrates that all-sky emissivity and DLR can be effectively estimated using shortwave measurements alone, significantly broadening the accessibility of longwave radiation modeling to locations lacking comprehensive instrumentation.

The proposed approach offers promising implications for enhancing climate simulations, solar forecasting, and remote sensing applications, especially in data-sparse areas. Future work may focus on integrating satellite-derived cloud properties or testing hybrid parameterization strategies to further improve predictive accuracy.

\section*{Data Availability}

The radiation dataset used in this study is publicly archived on the open-source data repository Zenodo 
and will be made available following peer review acceptance.

\noindent DOI: \href{https://doi.org/10.5281/zenodo.xxxxxx}{10.5281/zenodo.xxxxxx}

\noindent The dataset comprises 4,905,673 observations of simultaneous global horizontal 
irradiance (GHI), diffuse horizontal irradiance (DHI), direct normal irradiance (DNI), 
and downwelling longwave radiation (DLR) measurements from seven SURFRAD stations 
(Bondville, Desert Rock, Fort Peck, Goodwin Creek, Penn State University, Sioux Falls, 
and Table Mountain) across diverse climate zones spanning 2010--2024. Data are provided 
in HDF5 format with compression (total size: 15.47 GB). Complete documentation including 
data dictionary, quality control procedures, metadata following CF Conventions, and 
processing scripts are provided with the dataset.

The dataset is currently under embargo until [EMBARGO DATE, e.g., July 31, 2026] to protect 
confidentiality during the peer review. Reviewers and editors may access the data via a private 
link provided to the journal. Upon manuscript acceptance, the embargo will be automatically 
lifted, and the data will be freely available to the research community in perpetuity. 
The dataset is licensed under Creative Commons Attribution 4.0 (CC-BY-4.0), permitting unrestricted use, distribution, and reproduction, provided appropriate attribution is given.

Data loading and analysis routines in Python and MATLAB are provided in the repository 
to facilitate reproducibility and ease of use \citep{SURFRAD_data_cleaned}. Furthermore, all the scripts utilized for data processing, model fitting, validation and plotting is made available in a public GitHub repository at \url{NEEDTOADDURL}.

\section*{CRediT Author Contribution Statement}

\textbf{Akarsh Srivastava:} Need to write
\textbf{Jiedong Wang:} Need to write
Akarsh Srivastava and Jiedong Wang contributed equally to this work and share first authorship.
\textbf{Zoey Huang:} Need to write
\textbf{Clayson Briggs:} Need to write
\textbf{Carlos F.\,M. Coimbra:} Need to write

%% Appedix %%
\newpage

\appendix
\section{DRA data availability}
Station DRA was excluded from the model validation since its station-specific fit significantly deviated from the overall fit, largely due to the lack of cloud cover and corresponding cloudy-sky data at the station. To further quantify the lack of cloudy-sky data, data from each station are sorted into bins corresponding to $k_d$ and shown in Fig.~\ref{fig:buckets}.

As seen in Fig.~\ref{fig:buckets}, most stations except for DRA contribute relatively equally to the data bins. In contrast, DRA contributes a significant amount of data to the 0-0.2 bin and vanishingly little data to each bin with $k_d$ values $>0.2$. These low values of $k_d$ mean that the shortwave radiation incident on DRA usually had a small diffuse component, indicating mostly clear skies. This corroborates the decision to exclude DRA from the validation analysis, since the station-specific DRA model is largely only valid for data with $k_d$ values between 0 and 0.2. 

\begin{figure}[!htbp]
    \centering
    \includegraphics[width=1\textwidth, keepaspectratio]{paper_pic/kd_buckets.png}
    \caption{Total counts of data used in the overall model fit by station and $k_d$ buckets. Values of $k_d$ close to 0 correspond to low amounts of diffuse shortwave radiation and correlate with less cloud cover, while values of $k_d$ close to 1 correspond to high amounts of diffuse shortwave radiation and correlate with more cloud cover.}
    \label{fig:buckets}
\end{figure}

Furthermore, parameters for the radiative model were found again using two different bucketing methods. The first bucketing method used weighted buckets, finding parameters of $c_1 = 0.603$ and $c2 = 1.922$ resulting in $\mathrm{RMSE} = 3.99 \%$. The second method used equal data points per bucket, finding $c_1 = 0.608$ and $c_2 = 1.865$ for a result of $\mathrm{RMSE} = 5.01 \%$. Both methods result in low errors and the parameters found are similar to those reported in the main body, confirming that the inclusion of DRA data and the binning methods used do not significantly affect the model results. 

\section{Extended error analysis using training data}
To thoroughly evaluate the applicability, accuracy, and stability of the longwave radiation prediction model proposed in this study across various spatial locations and time scales, an in-depth error analysis was conducted using the model proposed in Eq.~\eqref{final_model}. Based on annual data from 2010 to 2022 collected from the seven representative SURFRAD stations, the relative mean bias error (rMBE) and the relative root mean square error (rRMSE) were systematically employed to assess the differences between model predictions and observed measurements. To more clearly illustrate the temporal and spatial variability and trends in error distributions, the results are presented with heatmaps in Fig.~\ref{fig:error_heatmaps}.

\begin{figure}[!htbp]% placement specifier
%% Use \includegraphics command to insert graphic files. Place graphics files in 
%% working directory.
\centering%% For centre alignment of image.
\includegraphics[width=1\textwidth, keepaspectratio]{paper_pic/heatmap_rMBE_rRMSE_kd_model.png}
%% Use \caption command for figure caption and label.
  \caption{
    Heatmaps showing annual prediction errors of the proposed longwave radiation model at the seven SURFRAD stations from 2010 to 2022 with two error metrics: 
    (a) relative mean bias error (rMBE) and 
    (b) relative root mean squared error (rRMSE). 
    Note that data for the DRA station are partially missing between 2017 and 2022 according to the discussion in Section 3.2.1 as indicated by the blank areas in the heatmaps.
  }
  \label{fig:error_heatmaps}
\end{figure}

The rMBE values shown in Fig.~\ref{fig:error_heatmaps}(a) for most stations and years consistently remain within ±1.0, indicating that the model generally exhibits minimal systematic bias and closely aligns with observed longwave radiation. Notably, the BON and SXF stations maintain near zero bias values throughout most years, highlighting the model's exceptional accuracy and stable performance in these regions. However, certain stations such as TBL and FPK experience more pronounced bias fluctuations in specific years, potentially attributable to local weather events, anomalies in microclimates, or measurement instrument errors.

The rRMSE values in Fig.~\ref{fig:error_heatmaps}(b) predominantly range between 3.0 and 4.5, indicating stable and efficient predictive capability. Particularly noteworthy is that stations such as BON, GWN, and DRA (with complete data before 2017) consistently exhibit low and stable rRMSE values, underscoring the model's robustness and generalizability across diverse climatic and observational environments. Even during complex meteorological conditions and years with abnormal climatic events (e.g., 2016 and 2017), the model did not display significant performance degradation, demonstrating its capacity to effectively handle extreme or varying climate scenarios.

In conclusion, the error analysis results presented in this section show low systematic biases and minimal error fluctuations across most stations and years, further validating the reliability and predictive power of the proposed radiative model. 

% \section{Worked Example: DLR Calculation}

% We provide a worked example below to demonstrate the practical application of the proposed model for estimating downwelling longwave radiation (DLR) under typical summer conditions at the Bondville, IL (BON) station. This example uses typical summer conditions at BON, including air temperature, relative humidity, site altitude, and shortwave radiation measurements, to compute the DLR step-by-step using the proposed model. It aids users in implementing the model and verifying their calculations and the work flow outlined in Algorithm~\ref{alg:dlr_estimation}.

% \begin{tcolorbox}[
%     colback=gray!5!white,
%     colframe=black!60,
%     breakable,
%     enhanced,
%     sharp corners,
%     boxrule=0.4pt,
%     left=6pt,
%     right=6pt,
%     top=6pt,
%     bottom=6pt
% ]

% \textbf{Given (Bondville, IL -- typical summer conditions):}
% \begin{itemize}
%     \item $T_a = 298.15$ K (25$^\circ$C)
%     \item $\phi = 0.60$ (60\% relative humidity)
%     \item $z = 213$ m (site altitude)
%     \item GHI $= 650$ W/m$^2$, DHI $= 195$ W/m$^2$
% \end{itemize}

% \textbf{Solution:}

% \textit{Step 1--2: Water vapor pressure}
% \begin{align*}
% P_s &= 610.94 \exp\left(\frac{17.625 \times 25}{298.15 - 30.11}\right) = 3168.7 \text{ Pa} \\
% P_w &= 3168.7 \times 0.60 = 1901.2 \text{ Pa} \\
% p_w &= 1901.2 / 101325 = 0.01877
% \end{align*}

% \textit{Step 3: Clear-sky emissivity}
% \begin{align*}
% \varepsilon_{\text{clear}} &= 0.600 + 1.652\sqrt{0.01877} + 0.150\left[\exp\left(-\frac{213}{8500}\right) - 1\right] \\
% &= 0.600 + 0.226 - 0.0037 = \mathbf{0.822}
% \end{align*}

% \textit{Step 4: Diffuse fraction}
% \[
% k_d = 195 / 650 = \mathbf{0.30}
% \]

% \textit{Step 5: Cloud fraction factor}
% \[
% \gamma = 0.585 \times 0.30^{1.748} = 0.585 \times 0.123 = \mathbf{0.072}
% \]

% \textit{Step 6: All-sky emissivity}
% \[
% \varepsilon_{\text{sky}} = 0.822 + 0.072 \times (1 - 0.822) = 0.822 + 0.013 = \mathbf{0.835}
% \]

% \textit{Step 7: DLR}
% \[
% \text{DLR} = 0.835 \times 5.67 \times 10^{-8} \times 298.15^4 = \mathbf{374.4 \text{ W/m}^2}
% \]

% \end{tcolorbox}

% include bibliography
\bibliographystyle{elsarticle-num} 
\bibliography{reference}

\end{document}

\endinput
%%
%% End of file `elsarticle-template-harv.tex'.